\documentclass{beamer}

% \usetheme{Pittsburgh}
 \setbeamertemplate{enumerate items}[circle]
%  \usetheme{Boadilla}
\usepackage[ngerman,english]{babel}			
\usepackage[utf8]{inputenc}				

\usepackage{color}
\usepackage{booktabs}
\usepackage{animate}
\usepackage{hyperref}
\usepackage{dirtytalk} % quotes
\usepackage{amsmath}

% \usepackage{natbib}
% \usepackage{apalike}
% \bibliography{refs}

\usepackage[style=apa]{biblatex}
\addbibresource{refs.bib}

% Define command to print full bibliography in footnote
\newcommand{\fcite}[1]{\footnote{\fullcite{#1}}}

\hypersetup{colorlinks,urlcolor=blue, linkcolor=blue}

\setbeamertemplate{footline}[frame number]
\setbeamertemplate{section in toc}[sections numbered]
\setcounter{tocdepth}{1}
% \setbeamertemplate{subsection in toc}[subsections numbered]

% LOGO TOP RIGHT
% \usepackage{tikz}
% \titlegraphic { 
% \begin{tikzpicture}[overlay,remember picture]
% \node[left=0.2cm] at (current page.30){
%     
\includegraphics[width=3cm]{figures/mpidr_logo_colour_de_1.png}
% };
% \end{tikzpicture}
% }

% LOGO BOTTOM RIGHT
\titlegraphic{

\includegraphics[width=5cm]{figures/mpidr_logo_colour_de_1.png}
\hfill

\includegraphics[width=3cm]{figures/edsd-logo.png}
}

% \makeatletter
% \setbeamertemplate{navigation symbols}{}

% make blocks no transparent
\setbeamercolor{block body alerted}{bg=alerted text.fg!10}
\setbeamercolor{block title alerted}{bg=alerted text.fg!20}
\setbeamercolor{block body}{bg=structure!10}
\setbeamercolor{block title}{bg=structure!20}
\setbeamercolor{block body example}{bg=green!10}
\setbeamercolor{block title example}{bg=green!20}



%%	title page information

\title[Kinship]{Kinship Structures (1)}
\subtitle{Introduction to Kinship Demography}

\author{Diego Alburez Guti\'errez$^{\dagger}$}

\institute{$^{\dagger}$Kinship Inequalities Research Group, \\ Max Planck Institute for Demographic Research }

\date{European Doctoral School of Demography 2024-25 \\ 18 Nov 2024}

%%	Begin of the document:
%%	======================

\begin{document}

 \begin{frame}[t,plain]
 \titlepage
 \end{frame}
 
 \begin{frame}{Consider a  baby born in France in 1950...}
     \begin{enumerate}
         \item How old were her grandparents when she was born, on average?
         \item How many living children did she have on her 70th birthday?
         \item How many grandchildren?
     \end{enumerate}
 \end{frame}

 \begin{frame}{Agenda}
\tableofcontents
\end{frame}

\section{Course setup and evaluation}
\begin{frame}
    \begin{centering}
    \begin{beamercolorbox}[sep=12pt,center]{part title}
    \usebeamerfont{section title}\insertsection\par
    \end{beamercolorbox}
    \end{centering}
\end{frame}

\begin{frame}{Course overview}

\begin{enumerate}
    \item Mon - ``Introduction to kinship demography''
    \item Tue - ``The formal demography of kinship''
    \item Wed - ``Extensions of the kinship model''
    \item Thu - Assignment
\end{enumerate}

\vspace{1cm}

Course website: See syllabus!
    
\end{frame}

\begin{frame}{Schedule}

\begin{itemize}
    \item Meet every day 9:00-12:00 
    \item Mon-Wed
    \begin{itemize}
        \item 1.5h lectures
        \item 1.5h lab session
    \end{itemize}
    \item Thursday (assignment work only)
\end{itemize}

\end{frame}

\begin{frame}{Lab session}
       \begin{figure}[!h]
       \centering
\includegraphics[width =1\textwidth]{figures/pack.JPG}
     \end{figure}
\end{frame}


\begin{frame}{Final assignment}

\textbf{Goal}

    \begin{itemize}
        \item Hands-on experience with kinship models 
        \item Reflect on model specifications, assumptions, and outcomes
    \end{itemize}
    
    \vspace{0.5cm}
    \textbf{Exercises}
    
    \begin{itemize}
        \item Model kinship dynamics using `DemoKin` package
        \item Benchmark models of kinship
    \end{itemize}
    
    \vspace{0.5cm}
    \textbf{See full instructions in syllabus!} 
    
\end{frame}

\section{Introduction to kinship demography}
\begin{frame}
    \begin{centering}
    \begin{beamercolorbox}[sep=12pt,center]{part title}
    \usebeamerfont{section title}\insertsection\par
    \end{beamercolorbox}
    \end{centering}
\end{frame}


\begin{frame}{Why do we need a course on kinship demography?}

Kinship demography, in recent years:

  \vspace{0.5cm}
  
   \begin{enumerate}
       \item theoretical, methodological, and substantive developments
       \item opportunity to advance discipline as a whole
       \item low visibility of kinship demography
   \end{enumerate}
\end{frame}

\begin{frame}{Great societal interest\footnote{`The age of the grandparent has arrived.'  (Jan 2023). The Economist. \url{https://www.economist.com/international/2023/01/12/the-age-of-the-grandparent-has-arrived}}}
\begin{figure}
\centering
\begin{minipage}{.5\textwidth}
  \centering
  \includegraphics[width=1\linewidth]{figures/econ2.JPG}
  \label{fig:test1}
\end{minipage}%
\begin{minipage}{.5\textwidth}
  \centering
  \includegraphics[width=1\linewidth]{figures/econ3.JPG}
  \label{fig:test2}
\end{minipage}
\end{figure}
\end{frame}

\begin{frame}{Definitions (1)\fcite{alburez-gutierrez_kinship_2022}}
   \begin{exampleblock}{Kinship}
        Social relationships that bind individuals together through culturally shared definitions of relatedness on biological, legal, or normative grounds, ultimately constituting family systems.
   \end{exampleblock}
   
      \vspace{0.5cm}
   
    \begin{alertblock}{Family}
        More narrow group of kin given special privilege which, among other things, organize the provision of support, socialization, and social placement of its members.
  \end{alertblock}

\end{frame}

\begin{frame}{The role of kinship in human societies}
   \begin{enumerate}
       \item Socialisation, protection, and sustenance
       \item Inter-generational solidarity: exchanges and bequests
       \item Social structure and identity
       \item Early-life conditions $\rightarrow$ later-life outcomes
   \end{enumerate}
\end{frame}



\begin{frame}{Definitions (2)\fcite{alburez-gutierrez_kinship_2022}}
   \begin{exampleblock}{Kinship demography}
        The study of family networks, their structures and dynamics from a demographic perspective and using demographic methods.
   \end{exampleblock}
\end{frame}

\begin{frame}{Kinship as a demographic human universal}
    \begin{enumerate}
        \item All humans are born 
        \item All humans die
        \item All humans are embedded in kinship structures\fcite{caswell_formal_2019}
        \item No particular family configuration is universal or stable
    \end{enumerate}
\end{frame}


% \begin{frame}{}
% % \flushright \fcite{alburez-gutierrez_kinship_2022}     
%      \begin{figure}[!h]
%        \centering
%        \includegraphics[width =1\textwidth]{figures/kinship_diag.pdf} 
%      \end{figure}
% \end{frame}

% \begin{frame}{Examples: Kinequalities by gender and race in the US}
% \begin{enumerate}
%     \item Disparities in grandparental availability\fcite{margolis_cohort_2019}
%     \item Prevalence of unemployment in family networks\fcite{song2021role}
%     \item Effect of kin death on education\fcite{patterson_linked_2020}
% \end{enumerate}
% \end{frame}


% \subsection{Fundamentals}
% \begin{frame}
%     \subsectionpage
% \end{frame}


\begin{frame}{The \textit{Strong Ergodic Principle}\fcite{sharpe_problem_1911}}

    \say{A closed population with unchanging mortality and fertility rates has an \textbf{implicit age structure}.}
    % intrinsic rate of natural increase $r$ and 
    \vspace{0.3cm}
    
    The proportion of the population aged $x$ to $x+n$ in a stable population is\fcite{wachter_essential_2014}:
    
    \begin{equation}
      \frac{_{n}K_{x}}{_{\infty}K_{0}} = b\frac{_{n}L_{x}}{l_{0}}e^{-rx}
    \end{equation}

\end{frame}

\begin{frame}{Implied age structures using Leslie matrices}
      \begin{figure}[!h]
      \centering
      \includegraphics[width =1\textwidth]{figures/leslie_pop_share.pdf}
     \end{figure}    
\end{frame}

\begin{frame}{The demographic foundations of kin structure}
       \begin{figure}[!h]
       \centering
       \includegraphics[width =1\textwidth]{figures/keyf.JPG}
     \end{figure}
\end{frame}


\begin{frame}{Implied kinship structures}
    
\vspace{0.3cm}
\say{A fixed set of age-specific rates implies the probability that a girl aged $a$ has a living mother and great-grandmother, as well as her expected number of daughters, sisters, aunts, nieces, and cousins.}\fcite{keyfitz_applied_1985}
\end{frame}


\begin{frame}{Expected number of relatives (USA 1965, F, 45 yo)\fcite{keyfitz_applied_1985}}
       \begin{figure}[!h]
       \centering
       \includegraphics[width =.75\textwidth]{figures/tree.JPG}
     \end{figure}
\end{frame}


% \begin{frame}{Expected kin count for Swedish woman born in...\footnote{\url{http://alburez.me/DemoKin_example/}}}
%       \begin{figure}[!h]
%       \centering
%       \animategraphics[width = 1\textwidth]{40}{figures/kin_count_anim}{0}{40}
%      \end{figure}
% \end{frame}

\begin{frame}{Expected kin count for a Swedish woman\footnote{Time-invariant conditions;  \url{http://alburez.me/DemoKin_example/}}}
       \begin{figure}[!h]
       \centering
       \includegraphics[width = 1\textwidth]{figures/kin_count_anim}
     \end{figure}
\end{frame}

\begin{frame}{Observed kin count from Swedish registers\fcite{kolk_swedish_2021}}
       \begin{figure}[!h]
       \centering
       \includegraphics[width =1\textwidth]{figures/stock.JPG}
     \end{figure}
\end{frame}

\begin{frame}{Summary}
    \begin{itemize}
        \item Kinship matters for individuals and for societies
        \item Kinship is a demographic universal 
        \item Implicit kinship structures in populations
    \end{itemize}
\end{frame}

\begin{frame}{\texttt{DemoKin}: matrix kinship models in R\fcite{williams_demokin_2023}}

     \begin{minipage}{0.7\textwidth}
        \begin{itemize}
            \item Time-(in)variant models
            \item One/two-sex models
            \item Multistate models
            \item Kin loss by cause of death
            % \item \url{https://github.com/IvanWilli/DemoKin}
            \item More in the lab session...
        \end{itemize}
    \end{minipage} \hfill
    \begin{minipage}{0.25\textwidth}
    \begin{figure}[!h]
       \flushright
       
\includegraphics[width =1\textwidth]{figures/DemoKin-Logo.png}
     \end{figure}
    \end{minipage}
\end{frame}

% -----------------------------------------------

\begin{frame}{Discuss}
    \begin{enumerate}
        \item What is `kinship demography' and how is it different from `family demography?'
    \end{enumerate}
\end{frame}

\begin{frame}{}
    \centering
    \huge Break
\end{frame}

\section{Example: The lived experience of kin loss, or looking at mortality through a kinship lens}
\begin{frame}
    \begin{centering}
    \begin{beamercolorbox}[sep=12pt,center]{part title}
    \usebeamerfont{section title}\insertsection\par
    \end{beamercolorbox}
    \end{centering}
\end{frame}

\begin{frame}{Death and loss}
    
\say{People living through [the Covid-19 pandemic] do not experience the world by weekly rates; they perceive mortality risk through the experiences of their own social network.}\fcite{trinitapoli_demography_2021}
    
\end{frame}

\begin{frame}{Deaths experienced at each age (Sweden 1950, F)\footnote{\url{http://alburez.me/DemoKin_example/}}}
       \begin{figure}[!h]
       \centering
       \includegraphics[width =1\textwidth]{figures/deaths.jpg}
     \end{figure}
\end{frame}


% -----------------------------------------------
\subsection{Mortality $\Rightarrow$ kin loss}
\begin{frame}
    \begin{centering}
    \begin{beamercolorbox}[sep=12pt,center]{part title}
    \usebeamerfont{section title}\insertsubsection\par
    \end{beamercolorbox}
    \end{centering}
\end{frame}

\begin{frame}{}
       \begin{figure}[!h]
       \centering
       \includegraphics[width =1\textwidth]{figures/pub_dem.JPG}
     \end{figure}
\end{frame}

\begin{frame}{Experience of child loss\fcite{alburez-gutierrez_womens_nodate}}

\begin{exampleblock}{Research question}
How can we characterise the experience of offspring loss from the vantage point of mothers?
\end{exampleblock}

 \begin{enumerate}
     \item Methods: formal methods from kinship demography
     \item Data: 2019 UN World Population Prospects estimates (1950-2020) and projections (2020-2100) for 201 countries 
    %  \item Results: prevalence and incidence of offspring loss over life, across cohorts, and around the world
 \end{enumerate}
\end{frame}

\begin{frame}{Demographic foundations of lived experience of kin loss}

    \textbf{Child Death} $CD_{(a,c)}$ is the cumulative number of child deaths experienced by a woman in cohort $c$, conditional on her survival to age \textit{a}:
    
    \begin{equation*}
    Children\:born_{(a,c)} = Children\:surviving_{(a,c)} + Children\:died_{(a,c)}
    \end{equation*}
  
\end{frame}


\begin{frame}{Demographic foundations of lived experience of kin loss}

\textbf{Child Death} $CD_{(a,c)}$ is the cumulative number of child deaths experienced by a woman in cohort $c$, conditional on her survival to age \textit{a}:

\begin{equation*}
Children\:died_{(a,c)} = Children\:born_{(a,c)} - Children\:surviving_{(a,c)}
\end{equation*}
  
\end{frame}


% \begin{frame}{The Child Death identity}
\begin{frame}{Demographic foundations of lived experience of kin loss}

\textbf{Child Death} $CD_{(a,c)}$ is the cumulative number of child deaths experienced by a woman in cohort $c$, conditional on her survival to age \textit{a}\fcite{goodman_family_1974}:

\begin{equation*}
\label{eq:CD}
\underbrace{CD_{(a,c)}}_{\text{Children died}} = \underbrace{\sum_{x=15}^{x=a} {_1F_{(x,c)}}}_{\text{Children born}}-\underbrace{\sum_{x=15}^{x=a} {_1F_{(x,c)}} l_{(a-x,c+x)}}_{\text{Children surviving}}
\end{equation*}
  
 where

\begin{itemize}
  \item $_1F_{(x,c)}$ are age-specific fertility rates for mothers
  \item $l_{(a-x,c+x)}$ are age-specific mortality probabilities for mother's children
\end{itemize}

\end{frame}


% \begin{frame}{Offspring loss and survival over the life course\footnote{\textcolor{blue}{Alburez-Gutierrez, D.}, et al. (2021) ``Women's experience of child death: A global demographic perspective.'' \textbf{Demography}. Advanced Publication.}}
% \begin{figure}[!h]
%     \centering
%     \includegraphics[width =0.9\textwidth]{figures/cd1.JPG}
% \end{figure}
% \end{frame}

\begin{frame}{Offspring loss over the life course\fcite{alburez-gutierrez_womens_nodate}}
  \begin{center}
       \animategraphics[width = 0.65\textwidth]{10}{figures/fig2ab_trans}{0}{50}
  \end{center}
\end{frame}


\begin{frame}{Offspring loss over the life course\fcite{alburez-gutierrez_womens_nodate}}
  \begin{center}
      \includegraphics[width = 1\textwidth]{figures/fig2ab}
  \end{center}
\end{frame}


\subsection{Kin loss $\Rightarrow$ bereavement}
\begin{frame}
    \begin{centering}
    \begin{beamercolorbox}[sep=12pt,center]{part title}
    \usebeamerfont{section title}\insertsubsection\par
    \end{beamercolorbox}
    \end{centering}
\end{frame}

\begin{frame}{}
       \begin{figure}[!h]
       \centering
       \includegraphics[width =1\textwidth]{figures/pub_bmj.JPG}
     \end{figure}
\end{frame}

\begin{frame}{How many mothers (\textperthousand) have ever lost a child?\fcite{smith-greenaway_global_2021}}
\begin{figure}[!h]
    \centering
    \includegraphics[width =1\textwidth]{figures/cd2.JPG}
\end{figure}
\end{frame}

\begin{frame}{}
       \begin{figure}[!h]
       \centering
       \includegraphics[width =1\textwidth]{figures/pub_pnas.JPG}
     \end{figure}
\end{frame}

\begin{frame}{Demographic change and the lived experience of kin loss}
   \begin{enumerate}
       \item Excess mortality from Covid-19
       \item Some deaths would not have happened in the absence of the pandemic
       \item Some people would not have lost a relative (i.e., `bereaved')
   \end{enumerate}
   
\end{frame}

\begin{frame}{The `bereavement multiplier' approach}

\say{\ldots every death from COVID-19 will leave approximately nine bereaved.}\fcite{verdery_tracking_2020}

      \begin{figure}[!h]
      \centering
      \includegraphics[width =.5\textwidth]{figures/fig1a.pdf}
     \end{figure}
    
\end{frame}

\begin{frame}{`Excess bereavement': the survivor's perspective\fcite{snyder_estimates_2022}}
      \begin{figure}[!h]
       \centering
       \includegraphics[width =.5\textwidth]{figures/fig1b.pdf}
     \end{figure}
\end{frame}


% \begin{frame}{Kin loss in the context of Covid-19}
%       \begin{figure}[!h]
%       \centering
%       \includegraphics[width =1\textwidth]{figures/covid.png}
%      \end{figure}
% \end{frame}


\begin{frame}{How many people in the UK lost a relative to Covid-19?\fcite{snyder_estimates_2022}}
      \begin{figure}[!h]
       \centering
       \includegraphics[width =.9\textwidth]{figures/fig2_uk.pdf}
     \end{figure}
\end{frame}

% \begin{frame}{How many people lost a relative to Covid-19?\fcite{snyder_estimates_2022}}
% \centering
%     \begin{table}
%         \begin{tabular}{llr}
%         \toprule
%         Kin lost & Norway & Sweden \\
%         \midrule
%         Grandparents & .7 & 44\\
%         Parents & .3 & 20\\
%         Siblings & .1 & 10\\
%         Children & .1 & .8\\
%         \bottomrule
%         \end{tabular}
%     \caption{Bereaved individuals (1000's)}
%     \end{table}
%     % \begin{table}
%     %     \begin{tabular}{llr}
%     %     \toprule
%     %     Kin lost & Germany & UK \\
%     %     \midrule
%     %     Grandparents & 254 & 532\\
%     %     Parents & 154 & 261\\
%     %     Siblings & 825 & 143\\
%     %     Children & 9 & 20\\
%     %     \bottomrule
%     %     \end{tabular}
%     % \caption{Bereaved individuals (1000's)}
%     % \end{table}
% \end{frame}


% -------------------------------------------

% \subsection{Bereavement $\Rightarrow$ memory}
% \begin{frame}
%     \begin{centering}
%     \begin{beamercolorbox}[sep=12pt,center]{part title}
%     \usebeamerfont{section title}\insertsubsection\par
%     \end{beamercolorbox}
%     \end{centering}
% \end{frame}

% \begin{frame}{}
%       \begin{figure}[!h]
%       \centering
%       \includegraphics[width =1\textwidth]{figures/pub_dem_mem.JPG}
%      \end{figure}
% \end{frame}

% \begin{frame}{Survivors of the Rio Negro massacres\footnote{CIRMA Historical Archive, cited in Tecu Osorio (2012).}}
% \begin{figure}[!h]
%     \centering
%     \includegraphics[width =.8\textwidth]{figures/idp.png}
% \end{figure}
% \end{frame}

% \begin{frame}{Long-term consequences of genocide}
% % \fcite{alburez-gutierrez_diego_vivir_2021}
% \begin{figure}[!h]
%     \centering
%     \includegraphics[width =1\textwidth]{figures/pop_by_sex.png}
% \end{figure}
% \end{frame}


% \begin{frame}{Rio Negro inhabitants by relatedness to genocide victims\fcite{alburez-gutierrez_demographic_2022}}
% \begin{figure}[!h]
%     \centering
%     \includegraphics[width =1\textwidth]{figures/Fig4.pdf}
% \end{figure}
% \end{frame}

% \begin{frame}{Demographic drivers of historical memory\footnote{Image by Vivian Guzm\'an.}}
% \begin{figure}[!h]
%     \centering
%     \includegraphics[width =1\textwidth]{figures/pacoxom.jpg}
% \end{figure}
% \end{frame}

% \begin{frame}{Summary}
%  \begin{enumerate}
%     \item Kinship perspective: mortality $\rightarrow$ kin loss $\rightarrow$ breavement $\rightarrow$ memory
%      \item Mortality crises `accelerate' the experience of bereavement
%      \item Age-sex distribution of `excess bereavement' different from excess mortality
%      \item How are these losses experienced and how to support surviving kin?
% \end{enumerate}


% \begin{equation*}
%     \text{Excess Bereavement} \approx \text{Excess Mortality} \times \text{Kinship Structure}
% \end{equation*}

% \end{frame}

\begin{frame}{Mortality from a kinship perspective}
    \begin{enumerate}
        % \item Mortality and bereavement are different
        \item Mortality $\Rightarrow$ kin loss $\Rightarrow$ breavement $\Rightarrow$ memory
        \item Mortality crises `accelerate' the experience of bereavement
        \item Implications for individuals and populations
    \end{enumerate}
    
    \begin{equation*}
        \text{Kin Loss} \approx \text{Mortality} \times \text{Kinship Structure}
    \end{equation*}

\end{frame}


\begin{frame}{Discuss}
    \begin{enumerate}
        \item How could you apply a `kinship lens' to study other demographic processes?
        \item Think of 2-3 concrete examples
    \end{enumerate}
\end{frame}


% \begin{frame}{Concluding thoughts}
%   \begin{enumerate}
%       \item Demography will continue shaping kinship dynamics
%       \item amidst growing family complexity and geographic spread
%       \item demographers are ideally placed to study these changes
%   \end{enumerate}
% \end{frame}

\end{document}


